\section{Antipodal pairing}\label{app:cut}

Write $\theta=\pi-\phi$.  The arguments of the two shortest radial images are
$\pi-\phi$ and $-\pi-\phi$.  Their numerator is
\begin{align*}
&(\pi-\phi)e^{-(\pi-\phi)^2/(2t)}
-(\pi+\phi)e^{-(\pi+\phi)^2/(2t)}\\
&\quad=2e^{-(\pi^2+\phi^2)/(2t)}
\left[\pi\sinh\frac{\pi\phi}{t}
-\phi\cosh\frac{\pi\phi}{t}\right].
\end{align*}
The same computation gives \eqref{eq:cutpair} for every
$b_j=(2j+1)\pi$, $j\in\mathbb N_0$.

For $x\ge0$, comparison of derivatives proves
$\tanh x\ge x/(1+x)$.  Hence, for $\phi>0$,
\[
b\tanh\frac{b\phi}{t}
\ge\frac{b^2\phi}{t+b\phi}>\phi
\]
under condition $b^2>t+b\phi$.  The bracket in
\eqref{eq:cutpair} is therefore positive.  Dividing its Taylor expansion by
$\sin\phi$ gives the limit $b^2/t-1$, times the positive common exponential
factor.  All later odd pairs satisfy a stronger inequality than the principal
pair.

These are two terms of the affine-Weyl image sum, not two minimizing
geodesics.  At antipodal points of $S^3$, the initial tangent direction may
be any point of the unit $S^2$, so the minimizing geodesics form a
two-dimensional Morse--Bott family.  The exact limit and its differentiated
majorants are developed in Appendix~\ref{app:global-atlas}.
